Una substància radioactiva es desintegra segons l'equació següent (en el sistema internacional, SI):
\[ N = N_0\, e^{-0,0050t} \]

\begin{apartat}{1}
Expliqueu el significat de les magnituds que intervenen en aquesta equació i indiqueu el període de semidesintegració de la substància. Justifiqueu la resposta.
\end{apartat}

\begin{apartat}{1}
Si en un moment determinat la mostra conté $1,0\times 10^{28}$ nuclis d'aquesta substància, calculeu l'activitat que tindrà al cap de 4,0 hores.
\end{apartat}
